% !TeX root = ../main.tex
% Read docs/writing-guide.md for purpose, outline and common mistakes.
\chapter{Simulation}
\label{chap:simulation}

% Specify model assumptions, inputs and limits before showing outputs.
\section{Model and analysis method}
\label{sec:methodology}
The demonstration reduces each set of durations to its arithmetic mean.
Equation~\ref{eq:mean} adds the durations and divides the total by their count.
% equation supplies a number automatically; never type a number by hand.
\begin{equation}
\bar{t}=\frac{1}{n}\sum_{i=1}^{n}t_i .
\label{eq:mean}
\end{equation}
The observation $t_i$ and mean $\bar{t}$ are measured in seconds.
The count $n$ and index $i$ have no unit. With $n=3$, the sum for concept A
is \qty{180}{\second}; division by three gives \qty{60}{\second}.

\section{Fictional code example}
\label{sec:code-example}
Listing~\ref{lst:mean-example} implements only the arithmetic mean used in this
teaching example. It is independently written demonstration code and is unrelated
to any real thesis, company system or proprietary implementation.

% The file path is relative to main.tex. The caption creates the list entry.
% The label lets the surrounding prose reference the listing automatically.
\lstinputlisting[
  language=Python,
  caption={Fictional Python example that calculates an arithmetic mean.},
  label={lst:mean-example}
]{code/example_mean.py}

\section{Simulation boundary}
This arithmetic example is not a physical sand-flow simulation. A real simulation
chapter must define governing equations, boundary conditions, solver settings,
convergence checks and comparison with independent evidence.
